% !TeX program = lualatex
\documentclass[12pt,aspectratio=169]{beamer}
\usepackage[utf8]{inputenc}
\usepackage[T1]{fontenc}
\usepackage{amsmath,amsfonts,amssymb}
\usepackage{graphicx}
\usepackage{xcolor}
\usepackage{hyperref}
\usepackage{anyfontsize}
\usepackage{lmodern}
\usepackage{longtable}
\usepackage{array}
\usepackage{booktabs}

% ====== EXTREME FONT REDUCTION ======
\renewcommand{\tiny}{\fontsize{3.5}{4.5}\selectfont}
\renewcommand{\scriptsize}{\fontsize{4.5}{5.5}\selectfont}
\renewcommand{\footnotesize}{\fontsize{5.5}{6.5}\selectfont}
\renewcommand{\small}{\fontsize{6.5}{7.5}\selectfont}
\renewcommand{\normalsize}{\fontsize{7.5}{8.5}\selectfont}
\renewcommand{\large}{\fontsize{8.5}{9.5}\selectfont}
\renewcommand{\Large}{\fontsize{9.5}{10.5}\selectfont}
\renewcommand{\LARGE}{\fontsize{10.5}{12}\selectfont}
\renewcommand{\huge}{\fontsize{11.5}{13.5}\selectfont}
\renewcommand{\Huge}{\fontsize{13}{16}\selectfont}

% ====== COLORS ======
\definecolor{redcorrect}{RGB}{200,30,30}
\definecolor{bluecorrect}{RGB}{30,80,200}
\definecolor{greencheck}{RGB}{0,150,50}
\definecolor{lightgray}{RGB}{245,245,245}
\definecolor{darkgray}{RGB}{40,40,40}
\definecolor{softyellow}{RGB}{255,248,200}
\definecolor{deepblue}{RGB}{20,60,150}
\definecolor{darkred}{RGB}{150,20,20}
\definecolor{orangewarning}{RGB}{255,140,0}
\definecolor{correctgreen}{RGB}{0,120,40}
\definecolor{trapcolor}{RGB}{180,80,0}
\definecolor{purpleconcept}{RGB}{120,50,180}
\definecolor{tealhard}{RGB}{0,120,120}

\setbeamercolor{title}{fg=white,bg=redcorrect}
\setbeamercolor{frametitle}{fg=white,bg=redcorrect}
\setbeamercolor{block title}{fg=white,bg=bluecorrect}
\setbeamercolor{block body}{fg=darkgray,bg=lightgray}
\setbeamercolor{normal text}{fg=darkgray}
\usecolortheme{default}

% ====== CUSTOM COMMANDS ======
\newcommand{\correct}[1]{\textcolor{greencheck}{\textbf{\checkmark}} #1}
\newcommand{\keypoint}[1]{\textcolor{deepblue}{\textbf{#1}}}
\newcommand{\hardconcept}[1]{\textcolor{darkred}{\textbf{#1}}}
\newcommand{\examtraps}[1]{\textcolor{trapcolor}{\textbf{⚠ TRAP:}} #1}
\newcommand{\recallbox}[1]{\fcolorbox{deepblue}{blue!5}{\parbox{0.88\textwidth}{\centering \textbf{REMEMBER:} #1}}}
\newcommand{\stepbox}[1]{%
	\begin{center}
		\fcolorbox{darkgray}{white}{%
			\parbox{0.92\textwidth}{\vspace{0.03cm} #1 \vspace{0.03cm}}%
		}
	\end{center}
}
\newcommand{\answerbox}[1]{\fcolorbox{correctgreen}{green!8}{\parbox{0.88\textwidth}{\centering \textcolor{correctgreen}{\textbf{\checkmark\ CORRECT:}} #1}}}
\newcommand{\wronganswer}[1]{\textcolor{redcorrect}{\textbf{✘}} #1}
\newcommand{\trapbox}[1]{\fcolorbox{trapcolor}{orange!8}{\parbox{0.88\textwidth}{\centering \textcolor{trapcolor}{\textbf{⚠ EXAM TRAP:}} #1}}}
\newcommand{\keyrule}[1]{\textcolor{tealhard}{\textbf{🔑}} #1}

\begin{document}
	
	% ================================================================
	% SLIDE 1: TITLE
	% ================================================================
	\begin{frame}
		\centering
		\vspace{0.3cm}
		{\fontsize{18}{22}\selectfont \textbf{CAMS DOMAIN C}}\\[0.1cm]
		{\fontsize{14}{17}\selectfont \textbf{BUILDING AN AML/AFC COMPLIANCE PROGRAM}}\\[0.15cm]
		{\fontsize{9}{11}\selectfont \textit{ALL NEW QUESTIONS — COMPLETE EXPANDED TUTORIAL}}\\[0.08cm]
		\rule{4cm}{0.5pt}\\[0.08cm]
		{\fontsize{6}{7}\selectfont Data Extraction · Data Mining · MLAT · PPP · Data Integration · Internal/External Sources}\\[0.04cm]
		{\fontsize{6}{7}\selectfont Batch Screening · Risk-Based Authentication · National e-ID · Crypto Screening}\\[0.04cm]
		{\fontsize{6}{7}\selectfont AI Transition · FATF 6-Step · Financial Secrecy Index · EWRA · MLRO Role}\\[0.04cm]
		{\fontsize{6}{7}\selectfont Biometrics · Board Oversight · Terrorist Financing · Human Trafficking · TCSP}\\[0.04cm]
		{\fontsize{6}{7}\selectfont DNFBP Risk · Layering · Proactive Response · Structuring · Concentration Risk}
		\vspace{0.2cm}
	\end{frame}
	
	% ================================================================
	% SLIDE 2: DATA EXTRACTION — THE NORMALIZATION TRAP
	% ================================================================
	\begin{frame}
		\frametitle{\fontsize{11}{13}\selectfont \textbf{Data Extraction — The Normalization Trap}}
		\begin{block}{\fontsize{7}{9}\selectfont \textbf{The Hard Question — CAMS Exam Favorite}}
			\scriptsize
			Which of the following best describes a challenge during AFC data extraction?
		\end{block}
		\vspace{0.02cm}
		\answerbox{\large \textbf{Data may come in different formats, requiring normalization during integration.}}
		\vspace{0.02cm}
		\begin{block}{\fontsize{7}{9}\selectfont \textbf{Natural Explanation — Why You Got This Wrong}}
			\scriptsize
			\begin{itemize}
				\item Many candidates think data extraction is just \textbf{pulling data} from systems. But the \textbf{real challenge} is what happens \textbf{after} you pull it.
				\item Different systems use \textbf{different formats}: one system uses "US," another uses "USA," another uses "840" or "United States."
				\item \textbf{Normalization} is the process of making all this data \textbf{consistent} so it can be used together.
				\item Without normalization, your transaction monitoring system can't properly analyze transactions or identify patterns.
				\item \textbf{Key principle:} Data extraction is \textbf{not} the challenge — \textbf{data normalization} is.
				\item \textbf{Exam Trap:} The exam will offer "pulling data from systems" as a distractor. The real challenge is \textbf{normalization}.
			\end{itemize}
		\end{block}
		\vspace{0.02cm}
		\recallbox{\scriptsize \textbf{Key Rule:} Data extraction challenge = \textbf{different formats requiring normalization}.}
	\end{frame}
	
	% ================================================================
	% SLIDE 3: DATA MINING — THE EXTERNAL SOURCES TRAP
	% ================================================================
	\begin{frame}
		\frametitle{\fontsize{11}{13}\selectfont \textbf{Data Mining — The External Sources Trap}}
		\begin{block}{\fontsize{7}{9}\selectfont \textbf{The Hard Question — CAMS Exam Favorite}}
			\scriptsize
			Which are requirements for effective data mining and matching in AFC controls? (Select Two.)
		\end{block}
		\vspace{0.02cm}
		\answerbox{\large \textbf{1. Comparing extracted data against structured and unstructured risk sources. 2. Applying intelligent filters to reduce false positives during entity resolution.}}
		\vspace{0.02cm}
		\begin{block}{\fontsize{7}{9}\selectfont \textbf{Natural Explanation — Why You Got This Wrong}}
			\scriptsize
			\begin{itemize}
				\item \wronganswer{Focusing on internal customer files for mining purposes} — This is a \textbf{trap}. Effective data mining requires \textbf{both internal and external} sources.
				\item \textbf{Structured data:} Databases, spreadsheets, formatted fields.
				\item \textbf{Unstructured data:} Emails, documents, free-text fields, news articles.
				\item \textbf{Intelligent filters:} Reduce false positives during entity resolution (matching entities across different data sources).
				\item \textbf{Why external sources matter:} Internal data alone misses external risks like sanctions, PEPs, and adverse media.
				\item \textbf{Key principle:} Data mining is \textbf{broad}, not narrow. You need \textbf{external} data to catch \textbf{external} risks.
			\end{itemize}
		\end{block}
		\vspace{0.02cm}
		\recallbox{\scriptsize \textbf{Key Rule:} Data mining = \textbf{internal + external} data sources + \textbf{intelligent filters}.}
	\end{frame}
	
	% ================================================================
	% SLIDE 4: DATA INTEGRATION — THE ISOLATION TRAP
	% ================================================================
	\begin{frame}
		\frametitle{\fontsize{11}{13}\selectfont \textbf{Data Integration — The Isolation Trap}}
		\begin{block}{\fontsize{7}{9}\selectfont \textbf{The Hard Question — CAMS Exam Favorite}}
			\scriptsize
			Which help ensure effective integration of multiple data sources into compliance platforms? (Select Three.)
		\end{block}
		\vspace{0.02cm}
		\answerbox{\large \textbf{1. Maintaining a common data dictionary. 2. Mapping all incoming data to standardized field names. 3. Aligning control data requirements to system inputs.}}
		\vspace{0.02cm}
		\begin{block}{\fontsize{7}{9}\selectfont \textbf{Natural Explanation — Why You Got This Wrong}}
			\scriptsize
			\begin{itemize}
				\item \wronganswer{Separating internal and external data into isolated systems} — This is a \textbf{trap}. Separation \textbf{undermines} integration, it doesn't help it.
				\item \textbf{Common data dictionary:} Ensures all systems use the same definitions (e.g., "customer" means the same thing everywhere).
				\item \textbf{Field mapping:} Ensures data from source systems is translated correctly into target systems.
				\item \textbf{Aligning control requirements:} Ensures the monitoring system receives all necessary data.
				\item \textbf{Key principle:} Integration is about \textbf{connection}, not isolation.
				\item \textbf{Exam Trap:} "Isolating data" sounds secure but actually breaks integration.
			\end{itemize}
		\end{block}
		\vspace{0.02cm}
		\recallbox{\scriptsize \textbf{Key Rule:} Integration = \textbf{common dictionary + field mapping + aligned requirements}.}
	\end{frame}
	
	% ================================================================
	% SLIDE 5: INTERNAL vs EXTERNAL SOURCES — THE SCOPE TRAP
	% ================================================================
	\begin{frame}
		\frametitle{\fontsize{11}{13}\selectfont \textbf{Internal vs External Sources — The Scope Trap}}
		\begin{block}{\fontsize{7}{9}\selectfont \textbf{The Hard Question — CAMS Exam Favorite}}
			\scriptsize
			Which characteristics describe internal versus external sources used during AFC data extraction? (Select Two.)
		\end{block}
		\vspace{0.02cm}
		\answerbox{\large \textbf{1. Internal sources reflect observed customer behavior or system-held values. 2. External sources include sanction or adverse media lists from agencies.}}
		\vspace{0.02cm}
		\begin{block}{\fontsize{7}{9}\selectfont \textbf{Natural Explanation — Why You Got This Wrong}}
			\scriptsize
			\begin{itemize}
				\item \wronganswer{Internal data includes account balance snapshots, not user activity} — This is \textbf{too narrow}. Internal data includes \textbf{much more} than just account balances.
				\item \textbf{Internal sources:} Transaction histories, account balances, customer profiles, onboarding records, system-held values, observed customer behavior.
				\item \textbf{External sources:} Sanctions lists (OFAC, UN), PEP lists, adverse media databases, regulatory watchlists, FIU alerts.
				\item \textbf{Key principle:} Internal data provides \textbf{behavioral context}. External data provides \textbf{risk context}. Both are essential.
				\item \textbf{Exam Trap:} "Account balance snapshots only" is too narrow — internal data includes much more.
			\end{itemize}
		\end{block}
		\vspace{0.02cm}
		\recallbox{\scriptsize \textbf{Key Rule:} Internal = \textbf{behavioral context}. External = \textbf{risk context}. Both needed.}
	\end{frame}
	
	% ================================================================
	% SLIDE 6: MLAT — THE ADMISSIBILITY TRAP
	% ================================================================
	\begin{frame}
		\frametitle{\fontsize{11}{13}\selectfont \textbf{MLAT — The Admissibility Trap}}
		\begin{block}{\fontsize{7}{9}\selectfont \textbf{The Hard Question — CAMS Exam Favorite}}
			\scriptsize
			Two countries establish an MLAT. Country A has data protection rules preventing sharing of client identities. Country B bases investigations on BO disclosures. What challenge arises?
		\end{block}
		\vspace{0.02cm}
		\answerbox{\large \textbf{Evidence exchanged under the MLAT will NOT be admissible in Country B's legal system.}}
		\vspace{0.02cm}
		\begin{block}{\fontsize{7}{9}\selectfont \textbf{Natural Explanation — Why You Got This Wrong}}
			\scriptsize
			\begin{itemize}
				\item Most candidates think the problem is \textbf{Country A won't share the data}. That's \textbf{obvious}. The exam tests deeper understanding.
				\item The \textbf{real issue} is that even if Country A \textbf{does} share the data, Country B's courts may \textbf{reject} it.
				\item \textbf{Admissibility} is the legal standard for evidence in court.
				\item If the evidence doesn't meet Country B's evidentiary standards, it's \textbf{useless} — even if shared.
				\item \textbf{Key principle:} MLATs must address both \textbf{sharing} AND \textbf{admissibility}. The exam tests admissibility, not sharing.
				\item \textbf{Exam Trap:} "Country A's data protection rules will prevent sharing" is too obvious — the exam wants admissibility.
			\end{itemize}
		\end{block}
		\vspace{0.02cm}
		\recallbox{\scriptsize \textbf{Key Rule:} MLAT challenge = \textbf{admissibility}, not just sharing.}
	\end{frame}
	
	% ================================================================
	% SLIDE 7: PPP — THE REAL-TIME INTELLIGENCE TRAP
	% ================================================================
	\begin{frame}
		\frametitle{\fontsize{11}{13}\selectfont \textbf{PPP — The Real-Time Intelligence Trap}}
		\begin{block}{\fontsize{7}{9}\selectfont \textbf{The Hard Question — CAMS Exam Favorite}}
			\scriptsize
			An FIU wants to use PPP for real-time intelligence sharing to disrupt a terrorist financing network. Which approach is most aligned?
		\end{block}
		\vspace{0.02cm}
		\answerbox{\large \textbf{Establishing a formal platform with direct two-way communication on specific threats.}}
		\vspace{0.02cm}
		\begin{block}{\fontsize{7}{9}\selectfont \textbf{Natural Explanation — Why You Got This Wrong}}
			\scriptsize
			\begin{itemize}
				\item Many candidates choose \wronganswer{anonymized, aggregated format after threat is resolved} because it sounds \textbf{safer} or more \textbf{privacy-protecting}.
				\item But \textbf{that's not a PPP}. That's \textbf{after-the-fact reporting}.
				\item \textbf{PPP is about DISRUPTION} — stopping crimes \textbf{before} they happen.
				\item To disrupt, you need \textbf{real-time two-way communication}.
				\item Anonymized reporting \textbf{protects privacy} but \textbf{doesn't enable disruption}.
				\item \textbf{Key principle:} PPPs are about \textbf{real-time} intelligence sharing to \textbf{disrupt} active threats.
				\item \textbf{Exam Trap:} Anonymized reporting sounds responsible but misses the point of PPPs.
			\end{itemize}
		\end{block}
		\vspace{0.02cm}
		\recallbox{\scriptsize \textbf{Key Rule:} PPP = \textbf{real-time two-way communication} for disruption.}
	\end{frame}
	
	% ================================================================
	% SLIDE 8: CROSS-BORDER REGULATORY COORDINATION — THE JOINT REVIEW TRAP
	% ================================================================
	\begin{frame}
		\frametitle{\fontsize{11}{13}\selectfont \textbf{Cross-Border Coordination — The Joint Review Trap}}
		\begin{block}{\fontsize{7}{9}\selectfont \textbf{The Hard Question — CAMS Exam Favorite}}
			\scriptsize
			A cross-border financial institution is supervised by three different regulators across two regions. What integrated strategy should the regulators apply?
		\end{block}
		\vspace{0.02cm}
		\answerbox{\large \textbf{Regulators should coordinate to clarify scopes, align risk approaches, and execute joint reviews.}}
		\vspace{0.02cm}
		\begin{block}{\fontsize{7}{9}\selectfont \textbf{Natural Explanation — Why You Got This Wrong}}
			\scriptsize
			\begin{itemize}
				\item \wronganswer{Each regulator conducts independent reviews} — \textbf{Duplicative and inefficient}.
				\item \wronganswer{One regulator takes sole responsibility} — \textbf{Legally impossible}.
				\item \textbf{Cross-border supervision:} Governed by the \textbf{Basel Committee's Consolidated Supervision Framework}.
				\item \textbf{Joint reviews:} Coordinated examinations reduce burden and ensure consistency.
				\item \textbf{Key principle:} \textbf{Home-host supervision} ensures the lead regulator and host regulators coordinate effectively.
			\end{itemize}
		\end{block}
		\vspace{0.02cm}
		\recallbox{\scriptsize \textbf{Key Rule:} Cross-border = \textbf{coordinate + align + joint reviews}.}
	\end{frame}
	
	% ================================================================
	% SLIDE 9: BATCH SCREENING — THE VOLUME TRAP
	% ================================================================
	\begin{frame}
		\frametitle{\fontsize{11}{13}\selectfont \textbf{Batch Screening — The Volume Trap}}
		\begin{block}{\fontsize{7}{9}\selectfont \textbf{The Hard Question — CAMS Exam Favorite}}
			\scriptsize
			Which are true about batch screening? (Select Two.)
		\end{block}
		\vspace{0.02cm}
		\answerbox{\large \textbf{1. It helps reveal changing risks within the existing customer base. 2. It is designed to handle large volumes of data.}}
		\vspace{0.02cm}
		\begin{block}{\fontsize{7}{9}\selectfont \textbf{Natural Explanation — Why You Got This Wrong}}
			\scriptsize
			\begin{itemize}
				\item \wronganswer{It provides real-time transaction screening} — This is \textbf{real-time screening}, not batch screening.
				\item \textbf{Batch screening:} Processes data in bulk — typically the entire customer database — against sanctions lists, PEP lists.
				\item \textbf{Purpose:} Detects existing customers who have become risky (e.g., newly designated sanctions).
				\item \textbf{Volume-oriented:} Designed to handle millions of records.
				\item \textbf{Key principle:} Batch screening = \textbf{periodic re-screening} of the \textbf{entire customer base}.
				\item \textbf{Exam Trap:} "Real-time" is a distractor — batch is periodic, not real-time.
			\end{itemize}
		\end{block}
		\vspace{0.02cm}
		\recallbox{\scriptsize \textbf{Key Rule:} Batch screening = \textbf{large volumes + changing risks}. NOT real-time.}
	\end{frame}
	
	% ================================================================
	% SLIDE 10: RISK-BASED AUTHENTICATION — THE USER EXPERIENCE TRAP
	% ================================================================
	\begin{frame}
		\frametitle{\fontsize{11}{13}\selectfont \textbf{Risk-Based Authentication — The UX Trap}}
		\begin{block}{\fontsize{7}{9}\selectfont \textbf{The Hard Question — CAMS Exam Favorite}}
			\scriptsize
			How do risk-based authentication systems optimize user experience?
		\end{block}
		\vspace{0.02cm}
		\answerbox{\large \textbf{They prompt additional checks only when user behavior deviates from expectations.}}
		\vspace{0.02cm}
		\begin{block}{\fontsize{7}{9}\selectfont \textbf{Natural Explanation — Why You Got This Wrong}}
			\scriptsize
			\begin{itemize}
				\item \wronganswer{They eliminate all authentication requirements} — \textbf{NO.} Security still required.
				\item \wronganswer{They always require MFA for all users} — That's \textbf{not risk-based}, it's \textbf{uniform}.
				\item \textbf{Risk-based authentication (RBA)} uses \textbf{behavioral analytics} to establish a baseline of "normal" behavior.
				\item When behavior deviates, RBA prompts additional checks.
				\item Low-risk users have \textbf{frictionless} experiences; high-risk activities trigger additional authentication.
				\item \textbf{Key principle:} RBA balances \textbf{security and convenience}.
				\item \textbf{Exam Trap:} "Eliminate all authentication" is a trap — RBA still requires authentication.
			\end{itemize}
		\end{block}
		\vspace{0.02cm}
		\recallbox{\scriptsize \textbf{Key Rule:} RBA = \textbf{prompt only when behavior deviates}.}
	\end{frame}
	
	% ================================================================
	% SLIDE 11: NATIONAL e-ID — THE INTEGRATION TRAP
	% ================================================================
	\begin{frame}
		\frametitle{\fontsize{11}{13}\selectfont \textbf{National e-ID — The Integration Trap}}
		\begin{block}{\fontsize{7}{9}\selectfont \textbf{The Hard Question — CAMS Exam Favorite}}
			\scriptsize
			A fintech relies solely on device recognition for authentication. What additional measures could improve controls? (Select Two.)
		\end{block}
		\vspace{0.02cm}
		\answerbox{\large \textbf{1. Implementing document verification with biometric authentication. 2. Integrating with national electronic identity verification systems.}}
		\vspace{0.02cm}
		\begin{block}{\fontsize{7}{9}\selectfont \textbf{Natural Explanation — Why You Got This Wrong}}
			\scriptsize
			\begin{itemize}
				\item \wronganswer{Enhancing digital footprint analysis through IP verification} — IP verification is \textbf{already part} of device recognition.
				\item \wronganswer{Relying solely on social media verification} — Not a reliable identity verification method.
				\item \textbf{Document verification + biometrics:} Ensures identity authenticity.
				\item \textbf{National e-ID integration:} Leverages trusted government identity databases.
				\item \textbf{Key principle:} Identity verification = \textbf{document verification + biometrics + e-ID integration}.
				\item \textbf{Exam Trap:} IP verification is already part of device recognition — not an additional measure.
			\end{itemize}
		\end{block}
		\vspace{0.02cm}
		\recallbox{\scriptsize \textbf{Key Rule:} Identity verification = \textbf{document + biometrics + e-ID}.}
	\end{frame}
	
	% ================================================================
	% SLIDE 12: AI TRANSITION — THE DATA/PEOPLE/EXPLAINABILITY TRAP
	% ================================================================
	\begin{frame}
		\frametitle{\fontsize{11}{13}\selectfont \textbf{AI Transition — The Data/People/Explainability Trap}}
		\begin{block}{\fontsize{7}{9}\selectfont \textbf{The Hard Question — CAMS Exam Favorite}}
			\scriptsize
			Which actions are most important when transitioning to AI-based compliance tools? (Select Three.)
		\end{block}
		\vspace{0.02cm}
		\answerbox{\large \textbf{1. Ensuring data quality for training algorithms. 2. Managing change in employee workflows. 3. Verifying transparency and explainability in AI models.}}
		\vspace{0.02cm}
		\begin{block}{\fontsize{7}{9}\selectfont \textbf{Natural Explanation — Why You Got This Wrong}}
			\scriptsize
			\begin{itemize}
				\item \wronganswer{Selecting third-party global audit providers} — This is \textbf{less critical} than data, people, and explainability.
				\item \textbf{Data quality:} "Garbage in, garbage out" — AI is only as good as its training data.
				\item \textbf{Employee workflows:} AI changes how analysts work — must provide adequate training.
				\item \textbf{Explainability:} Regulators require \textbf{explainable AI} — must justify why an AI flagged something.
				\item \textbf{Key principle:} AI transition = \textbf{data + people + explainability}.
				\item \textbf{Exam Trap:} "Selecting audit providers" sounds important but is not the primary concern.
			\end{itemize}
		\end{block}
		\vspace{0.02cm}
		\recallbox{\scriptsize \textbf{Key Rule:} AI transition = \textbf{data + people + explainability}.}
	\end{frame}
	
	% ================================================================
	% SLIDE 13: FATF 6-STEP FRAMEWORK — THE COMPONENTS TRAP
	% ================================================================
	\begin{frame}
		\frametitle{\fontsize{11}{13}\selectfont \textbf{FATF 6-Step — The Component Trap}}
		\begin{block}{\fontsize{7}{9}\selectfont \textbf{The Hard Question — CAMS Exam Favorite}}
			\scriptsize
			Which are key components of FATF's six-step framework for national ML/TF risk assessments? (Select Three.)
		\end{block}
		\vspace{0.02cm}
		\answerbox{\large \textbf{1. Environmental scanning. 2. Horizon scanning. 3. Threat identification.}}
		\vspace{0.02cm}
		\begin{block}{\fontsize{7}{9}\selectfont \textbf{Natural Explanation — Why You Got This Wrong}}
			\scriptsize
			\begin{itemize}
				\item \wronganswer{Cross-border tax enforcement} — This \textbf{sounds plausible} because it's related to financial crime, but it's \textbf{not} a FATF risk assessment component.
				\item \wronganswer{Offshore asset repatriation} — Also sounds related, but \textbf{not} a component.
				\item \textbf{Environmental scanning:} Understanding the context (what's happening in the country?).
				\item \textbf{Horizon scanning:} Identifying emerging risks (what's coming?).
				\item \textbf{Threat identification:} Identifying specific threats (who are the bad actors?).
				\item \textbf{Key principle:} FATF's framework is about \textbf{understanding risk}, not enforcement actions.
				\item \textbf{Exam Trap:} "Cross-border tax enforcement" sounds related but is NOT part of the 6-step framework.
			\end{itemize}
		\end{block}
		\vspace{0.02cm}
		\recallbox{\scriptsize \textbf{Key Rule:} FATF 6-step = \textbf{environmental + horizon scanning + threat identification}.}
	\end{frame}
	
	% ================================================================
	% SLIDE 14: FINANCIAL SECRECY INDEX — THE FACTORS TRAP
	% ================================================================
	\begin{frame}
		\frametitle{\fontsize{11}{13}\selectfont \textbf{Financial Secrecy Index — The Factors Trap}}
		\begin{block}{\fontsize{7}{9}\selectfont \textbf{The Hard Question — CAMS Exam Favorite}}
			\scriptsize
			Which contribute to an elevated Financial Secrecy Index rating? (Select Three.)
		\end{block}
		\vspace{0.02cm}
		\answerbox{\large \textbf{1. Minimal BO reporting. 2. Extensive cross-border services to non-residents. 3. High secrecy-to-scale ratio.}}
		\vspace{0.02cm}
		\begin{block}{\fontsize{7}{9}\selectfont \textbf{Natural Explanation — Why You Got This Wrong}}
			\scriptsize
			\begin{itemize}
				\item \wronganswer{High ranking on global GDP without AML laws} — This \textbf{sounds} like it would matter, but GDP alone \textbf{doesn't create secrecy}.
				\item \textbf{Minimal BO reporting:} If you don't know who owns companies, that's secrecy.
				\item \textbf{Cross-border services to non-residents:} Providing banking services to people who don't live there enables secrecy.
				\item \textbf{High secrecy-to-scale ratio:} How much secrecy relative to financial activity — if a small jurisdiction has a lot of secrecy, that's a red flag.
				\item \textbf{Key principle:} FSI measures \textbf{actual secrecy}, not just size or wealth.
				\item \textbf{Exam Trap:} "High GDP" sounds important but doesn't create secrecy.
			\end{itemize}
		\end{block}
		\vspace{0.02cm}
		\recallbox{\scriptsize \textbf{Key Rule:} FSI factors = \textbf{BO reporting + cross-border services + secrecy-to-scale ratio}.}
	\end{frame}
	
	% ================================================================
	% SLIDE 15: EWRA — THE MLRO ROLE TRAP
	% ================================================================
	\begin{frame}
		\frametitle{\fontsize{11}{13}\selectfont \textbf{EWRA — The MLRO Role Trap}}
		\begin{block}{\fontsize{7}{9}\selectfont \textbf{The Hard Question — CAMS Exam Favorite}}
			\scriptsize
			Which function uses EWRA results to shape risk posture?
		\end{block}
		\vspace{0.02cm}
		\answerbox{\large \textbf{The Money Laundering Reporting Officer (MLRO).}}
		\vspace{0.02cm}
		\begin{block}{\fontsize{7}{9}\selectfont \textbf{Natural Explanation — Why You Got This Wrong}}
			\scriptsize
			\begin{itemize}
				\item Many candidates choose \wronganswer{Internal Audit Lead} because audit \textbf{sounds} like it would shape risk.
				\item But \textbf{Internal Audit} is the \textbf{third line of defense} — they assess controls, they don't shape risk posture.
				\item The \textbf{MLRO} is responsible for \textbf{managing AML risk} — they use EWRA to adjust controls, enhance monitoring, and allocate resources.
				\item \textbf{Key principle:} MLRO = \textbf{risk manager}. Internal Audit = \textbf{risk assessor}. Different roles.
				\item \textbf{Exam Trap:} Internal Audit sounds authoritative but is NOT the risk shaper.
			\end{itemize}
		\end{block}
		\vspace{0.02cm}
		\recallbox{\scriptsize \textbf{Key Rule:} MLRO = \textbf{shapes risk posture}. Internal Audit = \textbf{assesses} controls.}
	\end{frame}
	
	% ================================================================
	% SLIDE 16: TAX EVASION vs TAX AVOIDANCE — THE DISTINCTION TRAP
	% ================================================================
	\begin{frame}
		\frametitle{\fontsize{11}{13}\selectfont \textbf{Tax Evasion vs Tax Avoidance — The Distinction Trap}}
		\begin{block}{\fontsize{7}{9}\selectfont \textbf{The Hard Question — CAMS Exam Favorite}}
			\scriptsize
			A company failed to report offshore revenue and laundered funds through an unrelated third party. What conclusion is best supported?
		\end{block}
		\vspace{0.02cm}
		\answerbox{\large \textbf{The company committed tax evasion by concealing income.}}
		\vspace{0.02cm}
		\begin{block}{\fontsize{7}{9}\selectfont \textbf{Natural Explanation — Why You Got This Wrong}}
			\scriptsize
			\begin{itemize}
				\item Many candidates choose \wronganswer{legally avoided taxes through offshore structuring} because offshore structuring \textbf{sounds} like tax planning.
				\item But there's a \textbf{critical distinction}:
				\begin{itemize}
					\item \textbf{Tax avoidance:} Legal — you structure to minimize taxes \textbf{within the law}. Full disclosure required.
					\item \textbf{Tax evasion:} Illegal — you \textbf{conceal} income or assets. No disclosure.
				\end{itemize}
				\item \textbf{Failure to report + laundering through a third party} = concealment = \textbf{evasion}.
				\item \textbf{Key principle:} Tax avoidance requires \textbf{full disclosure}. Concealment = evasion = predicate offense for ML.
				\item \textbf{Exam Trap:} "Offshore structuring" sounds like avoidance but concealment makes it evasion.
			\end{itemize}
		\end{block}
		\vspace{0.02cm}
		\recallbox{\scriptsize \textbf{Key Rule:} Avoidance = legal with disclosure. Evasion = illegal with concealment.}
	\end{frame}
	
	% ================================================================
	% SLIDE 17: BIOMETRICS — THE COMPROMISED DATA TRAP
	% ================================================================
	\begin{frame}
		\frametitle{\fontsize{11}{13}\selectfont \textbf{Biometrics — The Compromised Data Trap}}
		\begin{block}{\fontsize{7}{9}\selectfont \textbf{The Hard Question — CAMS Exam Favorite}}
			\scriptsize
			Which risk is specifically associated with facial recognition compared to passwords?
		\end{block}
		\vspace{0.02cm}
		\answerbox{\large \textbf{Compromised biometric data is difficult to change or reset unlike passwords.}}
		\vspace{0.02cm}
		\begin{block}{\fontsize{7}{9}\selectfont \textbf{Natural Explanation — Why You Got This Wrong}}
			\scriptsize
			\begin{itemize}
				\item \wronganswer{The technology requires more storage capacity} — Not a unique risk.
				\item \wronganswer{Facial recognition is less secure} — Not necessarily true.
				\item \textbf{Passwords can be changed} — if compromised, you reset them.
				\item \textbf{Biometrics cannot be changed} — you can't get a new face or fingerprint.
				\item Once biometric data is compromised, it's \textbf{permanently compromised}.
				\item \textbf{Key principle:} Biometric risk = \textbf{irreplaceable} — can't change like passwords.
				\item \textbf{Exam Trap:} "Storage capacity" is a red herring — the real risk is irreplaceability.
			\end{itemize}
		\end{block}
		\vspace{0.02cm}
		\recallbox{\scriptsize \textbf{Key Rule:} Biometric risk = \textbf{can't change like passwords}.}
	\end{frame}
	
	% ================================================================
	% SLIDE 18: AUTHENTICATION TECHNOLOGIES — THE BIOMETRICS TRAP
	% ================================================================
	\begin{frame}
		\frametitle{\fontsize{11}{13}\selectfont \textbf{Authentication — The Biometrics Trap}}
		\begin{block}{\fontsize{7}{9}\selectfont \textbf{The Hard Question — CAMS Exam Favorite}}
			\scriptsize
			Which technologies enhance authentication and security in financial systems? (Select Three.)
		\end{block}
		\vspace{0.02cm}
		\answerbox{\large \textbf{1. Multi-factor authentication. 2. Physiological biometrics. 3. Behavioral biometrics.}}
		\vspace{0.02cm}
		\begin{block}{\fontsize{7}{9}\selectfont \textbf{Natural Explanation — Why You Got This Wrong}}
			\scriptsize
			\begin{itemize}
				\item \wronganswer{Password biometrics} — This is a \textbf{trap}. There is \textbf{no such thing} as "password biometrics." Passwords are \textbf{not biometrics}.
				\item \textbf{Multi-factor authentication (MFA)} = multiple factors (something you know, have, are).
				\item \textbf{Physiological biometrics} = physical traits (fingerprint, facial recognition, iris scan, voice).
				\item \textbf{Behavioral biometrics} = behavioral patterns (typing rhythm, mouse movements, gait).
				\item \textbf{Key principle:} Biometrics are \textbf{physical or behavioral}. Passwords are \textbf{knowledge-based}.
				\item \textbf{Exam Trap:} "Password biometrics" sounds plausible but does NOT exist.
			\end{itemize}
		\end{block}
		\vspace{0.02cm}
		\recallbox{\scriptsize \textbf{Key Rule:} MFA + physiological + behavioral = correct. "Password biometrics" = trap.}
	\end{frame}
	
	% ================================================================
	% SLIDE 19: BOARD OVERSIGHT — THE OPERATIONAL TRAP
	% ================================================================
	\begin{frame}
		\frametitle{\fontsize{11}{13}\selectfont \textbf{Board Oversight — The Operational Trap}}
		\begin{block}{\fontsize{7}{9}\selectfont \textbf{The Hard Question — CAMS Exam Favorite}}
			\scriptsize
			Which actions demonstrate effective board oversight? (Select Three.)
		\end{block}
		\vspace{0.02cm}
		\answerbox{\large \textbf{1. Receiving regular reports on risks. 2. Reviewing management's response to audit findings. 3. Challenging management on weak performance.}}
		\vspace{0.02cm}
		\begin{block}{\fontsize{7}{9}\selectfont \textbf{Natural Explanation — Why You Got This Wrong}}
			\scriptsize
			\begin{itemize}
				\item Many candidates choose \wronganswer{Participating in SAR investigations} because it sounds \textbf{engaged} and \textbf{involved}.
				\item But \textbf{boards provide OVERSIGHT}, not operational involvement.
				\item \textbf{Participating in SAR investigations} is \textbf{operational} — that's for compliance staff, not the board.
				\item \textbf{Board duties:}
				\begin{itemize}
					\item Receive reports (stay informed)
					\item Review audit responses (hold management accountable)
					\item Challenge management (ask tough questions)
				\end{itemize}
				\item \textbf{Key principle:} Board = \textbf{oversight and challenge}. NOT \textbf{doing the work}.
				\item \textbf{Exam Trap:} "Participating in SAR investigations" sounds engaged but is operational, not governance.
			\end{itemize}
		\end{block}
		\vspace{0.02cm}
		\recallbox{\scriptsize \textbf{Key Rule:} Board = \textbf{oversight}. NOT operational involvement.}
	\end{frame}
	
	% ================================================================
	% SLIDE 20: TERRORIST FINANCING — THE MOVEMENT METHODS TRAP
	% ================================================================
	\begin{frame}
		\frametitle{\fontsize{11}{13}\selectfont \textbf{Terrorist Financing — The Movement Methods Trap}}
		\begin{block}{\fontsize{7}{9}\selectfont \textbf{The Hard Question — CAMS Exam Favorite}}
			\scriptsize
			Which techniques are commonly used by terrorists to move or store money across borders avoiding detection? (Select Three.)
		\end{block}
		\vspace{0.02cm}
		\answerbox{\large \textbf{1. Fake charities. 2. Hawala (informal value transfer). 3. Physical cash transport across porous borders.}}
		\vspace{0.02cm}
		\begin{block}{\fontsize{7}{9}\selectfont \textbf{Natural Explanation — Why You Got This Wrong}}
			\scriptsize
			\begin{itemize}
				\item \wronganswer{Large wire transfers through regulated banks} — This is a \textbf{trap} because it \textbf{sounds} like a method, but it's \textbf{detectable}.
				\item Terrorists want to \textbf{avoid detection}, so they use \textbf{unregulated} methods.
				\item \textbf{Fake charities:} Disguise TF as humanitarian aid.
				\item \textbf{Hawala:} Informal value transfer outside regulated banking.
				\item \textbf{Physical cash transport:} Crosses borders avoiding detection.
				\item \textbf{Key principle:} Terrorists use \textbf{unregulated} and \textbf{hard-to-detect} methods.
				\item \textbf{Exam Trap:} "Large wire transfers" sounds like a method but terrorists avoid detection.
			\end{itemize}
		\end{block}
		\vspace{0.02cm}
		\recallbox{\scriptsize \textbf{Key Rule:} TF = \textbf{fake charities + Hawala + physical cash}. NOT regulated wire transfers.}
	\end{frame}
	
	% ================================================================
	% SLIDE 21: HUMAN TRAFFICKING — THE LAUNDERING METHODS TRAP
	% ================================================================
	\begin{frame}
		\frametitle{\fontsize{11}{13}\selectfont \textbf{Human Trafficking — The Laundering Methods Trap}}
		\begin{block}{\fontsize{7}{9}\selectfont \textbf{The Hard Question — CAMS Exam Favorite}}
			\scriptsize
			Which characteristics enable human trafficking operations to effectively launder illicit proceeds? (Select Two.)
		\end{block}
		\vspace{0.02cm}
		\answerbox{\large \textbf{1. Use of mobile payment and remittance systems. 2. Weak cross-border enforcement cooperation.}}
		\vspace{0.02cm}
		\begin{block}{\fontsize{7}{9}\selectfont \textbf{Natural Explanation — Why You Got This Wrong}}
			\scriptsize
			\begin{itemize}
				\item \wronganswer{High correlation between trafficking and counterfeit document operations} — This is \textbf{related} to trafficking, but \textbf{counterfeit documents} are a \textbf{separate crime}, not a laundering method.
				\item \textbf{Mobile payments/remittances:} Enable rapid movement of small amounts (structuring).
				\item \textbf{Weak cross-border enforcement:} Allows traffickers to operate across jurisdictions.
				\item \textbf{Key principle:} Traffickers launder proceeds through \textbf{rapid, small, cross-border} transactions.
				\item \textbf{Exam Trap:} "Counterfeit documents" is related but not a laundering method.
			\end{itemize}
		\end{block}
		\vspace{0.02cm}
		\recallbox{\scriptsize \textbf{Key Rule:} Trafficking laundering = \textbf{mobile payments} + \textbf{weak cross-border enforcement}.}
	\end{frame}
	
	% ================================================================
	% SLIDE 22: TCSP — THE OUTSOURCING TRAP
	% ================================================================
	\begin{frame}
		\frametitle{\fontsize{11}{13}\selectfont \textbf{TCSP — The Outsourcing Trap}}
		\begin{block}{\fontsize{7}{9}\selectfont \textbf{The Hard Question — CAMS Exam Favorite}}
			\scriptsize
			A TCSP is reviewing a client who wants confidentiality and is linked to high-risk industries in the NRA. What should the TCSP do?
		\end{block}
		\vspace{0.02cm}
		\answerbox{\large \textbf{Decline immediately unless full BO info and risk justifications are disclosed and verified.}}
		\vspace{0.02cm}
		\begin{block}{\fontsize{7}{9}\selectfont \textbf{Natural Explanation — Why You Got This Wrong}}
			\scriptsize
			\begin{itemize}
				\item \wronganswer{Refer to offshore local counsel and detach from compliance responsibility} — This is a \textbf{trap}. You \textbf{cannot} outsource AML responsibility.
				\item \textbf{TCSP retains compliance responsibility} — always.
				\item \textbf{Confidentiality requests + high-risk industries = red flags.}
				\item \textbf{FATF Rec 24/25:} BO identification is \textbf{mandatory}.
				\item \textbf{Key principle:} You can use local counsel, but you \textbf{cannot transfer liability}. You remain responsible.
				\item \textbf{Exam Trap:} "Refer to local counsel" sounds reasonable but doesn't transfer liability.
			\end{itemize}
		\end{block}
		\vspace{0.02cm}
		\recallbox{\scriptsize \textbf{Key Rule:} TCSP cannot outsource AML responsibility. Decline unless BO disclosed.}
	\end{frame}
	
	% ================================================================
	% SLIDE 23: DNFBP RISK MANAGEMENT — THE EDD TRAP
	% ================================================================
	\begin{frame}
		\frametitle{\fontsize{11}{13}\selectfont \textbf{DNFBP Risk Management — The EDD Trap}}
		\begin{block}{\fontsize{7}{9}\selectfont \textbf{The Hard Question — CAMS Exam Favorite}}
			\scriptsize
			Which strategies help manage DNFBP risks? (Select Two.)
		\end{block}
		\vspace{0.02cm}
		\answerbox{\large \textbf{1. Creating a reputational risk/client selection committee. 2. Using standardized sector-specific questionnaires.}}
		\vspace{0.02cm}
		\begin{block}{\fontsize{7}{9}\selectfont \textbf{Natural Explanation — Why You Got This Wrong}}
			\scriptsize
			\begin{itemize}
				\item \wronganswer{Defaulting to EDD for all DNFBPs} — This \textbf{sounds} prudent, but it's \textbf{not risk-based}. You would over-control low-risk DNFBPs.
				\item \textbf{Risk-based approach:} Apply proportionate measures based on actual risk.
				\item \textbf{Client selection committee:} Governance for high-risk decisions.
				\item \textbf{Standardized questionnaires:} Consistent risk assessment across DNFBPs.
				\item \textbf{Key principle:} Risk management = \textbf{proportionate measures}, not blanket EDD.
				\item \textbf{Exam Trap:} "Defaulting to EDD" sounds prudent but is not risk-based.
			\end{itemize}
		\end{block}
		\vspace{0.02cm}
		\recallbox{\scriptsize \textbf{Key Rule:} DNFBP risk = \textbf{committee + standardized questionnaires}. NOT blanket EDD.}
	\end{frame}
	
	% ================================================================
	% SLIDE 24: LAYERING — THE CLASSIFICATION TRAP
	% ================================================================
	\begin{frame}
		\frametitle{\fontsize{11}{13}\selectfont \textbf{Layering — The Classification Trap}}
		\begin{block}{\fontsize{7}{9}\selectfont \textbf{The Hard Question — CAMS Exam Favorite}}
			\scriptsize
			A private loan extended to a shell company with no operations, loan amount matches criminal proceeds. What is the initial classification?
		\end{block}
		\vspace{0.02cm}
		\answerbox{\large \textbf{A layering tactic designed to conceal the origin of illicit funds.}}
		\vspace{0.02cm}
		\begin{block}{\fontsize{7}{9}\selectfont \textbf{Natural Explanation — Why You Got This Wrong}}
			\scriptsize
			\begin{itemize}
				\item \wronganswer{A placement strategy to integrate clean capital} — This is a \textbf{trap}. Placement is the \textbf{first stage} — entering illicit funds into the system.
				\item \textbf{Layering} is the \textbf{second stage} — concealing the origin.
				\item A \textbf{private loan to a shell company} creates a \textbf{false paper trail} that appears legitimate.
				\item The loan \textbf{distances} the funds from their criminal source.
				\item \textbf{Key principle:} Layering = \textbf{concealing origin}. Placement = \textbf{entering system}. Integration = \textbf{using funds legitimately}.
				\item \textbf{Exam Trap:} "Placement" sounds plausible but the loan is layering, not placement.
			\end{itemize}
		\end{block}
		\vspace{0.02cm}
		\recallbox{\scriptsize \textbf{Key Rule:} Layering = \textbf{concealing origin}. Placement = \textbf{entering system}.}
	\end{frame}
	
	% ================================================================
	% SLIDE 25: PROACTIVE RESPONSE — THE UPDATING TRAP
	% ================================================================
	\begin{frame}
		\frametitle{\fontsize{11}{13}\selectfont \textbf{Proactive Response — The Updating Trap}}
		\begin{block}{\fontsize{7}{9}\selectfont \textbf{The Hard Question — CAMS Exam Favorite}}
			\scriptsize
			What is a proactive institutional response to newly detected AML risks?
		\end{block}
		\vspace{0.02cm}
		\answerbox{\large \textbf{Immediately updating the enterprise-wide risk assessment.}}
		\vspace{0.02cm}
		\begin{block}{\fontsize{7}{9}\selectfont \textbf{Natural Explanation — Why You Got This Wrong}}
			\scriptsize
			\begin{itemize}
				\item \wronganswer{Reducing reliance on transaction monitoring systems} — This is \textbf{reactive} and \textbf{counterproductive}.
				\item \textbf{Proactive} = take action \textbf{before} the risk materializes.
				\item \textbf{Updating EWRA} allows the institution to:
				\begin{itemize}
					\item Adjust controls
					\item Allocate resources to higher-risk areas
					\item Enhance monitoring before risks materialize
				\end{itemize}
				\item \textbf{Key principle:} Proactive = \textbf{update risk assessment}. Reactive = \textbf{respond after the fact}.
				\item \textbf{Exam Trap:} "Reducing reliance on monitoring" sounds like a response but is reactive and wrong.
			\end{itemize}
		\end{block}
		\vspace{0.02cm}
		\recallbox{\scriptsize \textbf{Key Rule:} Proactive = \textbf{update EWRA immediately}.}
	\end{frame}
	
	% ================================================================
	% SLIDE 26: FINAL EXAM TIPS — DOMAIN C COMPLETE
	% ================================================================
	\begin{frame}
		\frametitle{\fontsize{11}{13}\selectfont \textbf{Domain C — Complete Exam Traps Summary}}
		\scriptsize
		\begin{block}{\fontsize{7}{9}\selectfont \textbf{All New Questions — Natural Explanations}}
			\scriptsize
			\begin{enumerate}
				\item \textbf{Data Extraction:} Challenge = \textbf{normalization}, not just pulling data.
				\item \textbf{Data Mining:} Need \textbf{internal + external} data sources + intelligent filters.
				\item \textbf{Data Integration:} Need \textbf{connection}, not isolation. Common dictionary + mapping + alignment.
				\item \textbf{Internal vs External:} Internal = behavioral context. External = risk context.
				\item \textbf{MLAT:} Challenge = \textbf{admissibility}, not just sharing.
				\item \textbf{PPP:} = \textbf{real-time two-way communication}, not post-hoc reporting.
				\item \textbf{Cross-Border:} = \textbf{coordinate + align + joint reviews}.
				\item \textbf{Batch Screening:} = \textbf{large volumes + changing risks}. NOT real-time.
				\item \textbf{Risk-Based Auth:} = \textbf{prompt only when behavior deviates}.
				\item \textbf{National e-ID:} = \textbf{document + biometrics + e-ID integration}.
				\item \textbf{AI Transition:} = \textbf{data + people + explainability}.
				\item \textbf{FATF 6-Step:} = \textbf{environmental + horizon scanning + threat identification}.
				\item \textbf{FSI:} = \textbf{BO reporting + cross-border services + secrecy-to-scale ratio}.
				\item \textbf{EWRA:} MLRO \textbf{shapes risk}; Internal Audit \textbf{assesses} controls.
				\item \textbf{Tax:} Avoidance = \textbf{legal with disclosure}. Evasion = \textbf{illegal with concealment}.
				\item \textbf{Biometrics:} Risk = \textbf{can't change like passwords}.
				\item \textbf{Authentication:} MFA + physiological + behavioral = correct. "Password biometrics" = trap.
				\item \textbf{Board:} = \textbf{oversight}, NOT operational involvement.
				\item \textbf{TF:} = \textbf{fake charities + Hawala + physical cash}. NOT regulated wire transfers.
				\item \textbf{Trafficking:} = \textbf{mobile payments + weak cross-border enforcement}.
				\item \textbf{TCSP:} Cannot outsource AML responsibility.
				\item \textbf{DNFBP:} = \textbf{committee + questionnaires}. NOT blanket EDD.
				\item \textbf{Layering:} = \textbf{concealing origin}. Placement = \textbf{entering system}.
				\item \textbf{Proactive:} = \textbf{update EWRA immediately}.
			\end{enumerate}
		\end{block}
		\vspace{0.02cm}
		\recallbox{\scriptsize \textbf{Final Rule:} Understand the \textbf{why} behind each concept. The exam tests \textbf{application}, not memorization.}
	\end{frame}
	
	% ================================================================
	% END SLIDE
	% ================================================================
	\begin{frame}
		\centering
		\vspace{1.5cm}
		{\fontsize{18}{22}\selectfont \textbf{Domain C — Complete}}\\[0.2cm]
		{\fontsize{11}{13}\selectfont All New Questions — Full Expanded Tutorial}\\[0.1cm]
		\rule{4cm}{0.5pt}\\[0.1cm]
		{\fontsize{8}{10}\selectfont Data Extraction · Data Mining · MLAT · PPP · Data Integration}\\[0.05cm]
		{\fontsize{8}{10}\selectfont Batch Screening · Risk-Based Auth · National e-ID · AI Transition}\\[0.05cm]
		{\fontsize{8}{10}\selectfont FATF 6-Step · FSI · EWRA · MLRO · Tax Evasion · Biometrics}\\[0.05cm]
		{\fontsize{8}{10}\selectfont Board Oversight · TF · Trafficking · TCSP · DNFBP · Layering}\\[0.05cm]
		{\fontsize{8}{10}\selectfont Proactive Response · Structuring · Concentration Risk · Shell Companies}\\[0.1cm]
		\rule{4cm}{0.5pt}\\[0.1cm]
		{\fontsize{8}{10}\selectfont \textit{All traps explained. All concepts covered. Ready for the exam.}}
		\vspace{0.5cm}
	\end{frame}
	
\end{document}